\documentclass[12pt]{article}
\usepackage[UTF8]{inputenc}
\usepackage{thaisupp}
\usepackage{enumitem}
\begin{document}
\title{ของไหล}
\maketitle
\section*{In-class Questions}
\begin{enumerate}[label=\arabic*.]
\item ความหนาแน่น (Density) คืออะไร \hfill \textbf{\small (Main)}
\begin{enumerate}[label=)]
\item ก) น้ำหนักต่อปริมาตร
\item ข) มวลต่อปริมาตร
\item ค) น้ำหนักต่อมวล
\item ง) ปริมาตรต่อมวล
\end{enumerate}
\item ความดัน (Pressure) คำนวณจากสูตรใด \hfill \textbf{\small (Main)}
\begin{enumerate}[label=)]
\item ก) P = F/A
\item ข) P = A/F
\item ค) P = F × A
\item ง) P = mgh
\end{enumerate}
\item ความดันในของเหลวที่ระดับลึก h จากผิวหาได้จากสูตรใด \hfill \textbf{\small (Main)}
\begin{enumerate}[label=)]
\item ก) P = ρgh
\item ข) P = ρgh/V
\item ค) P = mgh
\item ง) P = ρVg
\end{enumerate}
\item หลักของอาร์คิมีดีส (Archimedes' Principle) ระบุว่าอะไร \hfill \textbf{\small (Main)}
\begin{enumerate}[label=)]
\item ก) วัตถุที่จมในของเหลวจะได้รับแรงลอยเท่ากับน้ำหนักของมัน
\item ข) วัตถุที่จมในของเหลวจะได้รับแรงลอยเท่ากับน้ำหนักของของเหลวที่มันแทนที่
\item ค) วัตถุลอยได้เมื่อมีความหนาแน่นมากกว่าของเหลว
\item ง) แรงลอยขึ้นกับปริมาตรของวัตถุเท่านั้น
\end{enumerate}
\item หลักของเบอร์นูลลี (Bernoulli's Principle) ระบุว่าอะไรเมื่อของไหลไหลในท่อ \hfill \textbf{\small (Main)}
\begin{enumerate}[label=)]
\item ก) P + ½ρv² + ρgh = คงที่
\item ข) P × A = คงที่
\item ค) F = ρAV²
\item ง) P = ρgh
\end{enumerate}
\item สมการความต่อเนื่อง (Continuity Equation) A₁v₁ = A₂v₂ บ่งบอกอะไร \hfill \textbf{\small (Main)}
\begin{enumerate}[label=)]
\item ก) มวลของของไหลคงที่ในหน่วยเวลา
\item ข) ปริมาตรของของไหลคงที่ในหนึ่งวินาที
\item ค) ความดันของของไหลคงที่
\item ง) ความหนาแน่นของของไหลคงที่
\end{enumerate}
\item หลักของปาสคาล (Pascal's Principle) ระบุว่าอะไร \hfill \textbf{\small (Main)}
\begin{enumerate}[label=)]
\item ก) ความดันที่ให้กับของเหลวจะถูกส่งต่อไปทุกทิศทางเท่ากัน
\item ข) ของเหลวไหลจากที่สูงไปที่ต่ำเสมอ
\item ค) แรงลอยเท่ากับน้ำหนักของของเหลวที่แทนที่
\item ง) ความดันเพิ่มขึ้นตามความลึกเสมอ
\end{enumerate}
\item นักประดาน้ำลงไปที่ความลึก 20 m ใต้น้ำ ความดันที่เขาแบกรับเป็นเท่าไร (น้ำ: ρ = 1,000 kg/m³, g = 10 m/s², ความดันบรรยากาศ ≈ 1 × 10⁵ Pa) \hfill \textbf{\small (Main)}
\begin{enumerate}[label=)]
\item ก) 1 × 10⁵ Pa
\item ข) 2 × 10⁵ Pa
\item ค) 3 × 10⁵ Pa
\item ง) 5 × 10⁵ Pa
\end{enumerate}
\item วัตถุมีปริมาตร 0.002 m³ ลอยในน้ำ (ρ = 1,000 kg/m³) จงหาแรงลอย (g = 10 m/s²) \hfill \textbf{\small (Main)}
\begin{enumerate}[label=)]
\item ก) 2 N
\item ข) 20 N
\item ค) 200 N
\item ง) 0.2 N
\end{enumerate}
\item ท่อน้ำพื้นที่หน้าตัด 0.01 m² มีน้ำไหลด้วยความเร็ว 5 m/s ท่อหดเหลือพื้นที่ 0.005 m² ความเร็วน้ำที่ส่วนแคบเป็นเท่าไร \hfill \textbf{\small (Main)}
\begin{enumerate}[label=)]
\item ก) 2.5 m/s
\item ข) 5 m/s
\item ค) 10 m/s
\item ง) 20 m/s
\end{enumerate}
\item วัตถุมีมวล 200 g และปริมาตร 250 cm³ มีความหนาแน่นเท่าไร \hfill \textbf{\small (Main)}
\begin{enumerate}[label=)]
\item ก) 0.8 g/cm³
\item ข) 1.25 g/cm³
\item ค) 0.5 g/cm³
\item ง) 8 g/cm³
\end{enumerate}
\item ตามหลักของเบอร์นูลลี เมื่อของไหลในท่อไหลเร็วขึ้น ความดันจะเป็นอย่างไร \hfill \textbf{\small (Main)}
\begin{enumerate}[label=)]
\item ก) ความดันเพิ่มขึ้น
\item ข) ความดันลดลง
\item ค) ความดันคงที่
\item ง) ขึ้นกับชนิดของของไหล
\end{enumerate}
\item ทำไมเรือเหล็กจึงลอยน้ำได้ทั้งที่เหล็กมีความหนาแน่นมากกว่าน้ำ \hfill \textbf{\small (Main)}
\begin{enumerate}[label=)]
\item ก) เพราะเรือมีน้ำหนักเบากว่าน้ำที่มันแทนที่
\item ข) เพราะเรือมีรูปร่างพองาม
\item ค) เพราะเรือมีปริมาตรมากพอให้แรงลอยมากกว่าน้ำหนัก
\item ง) เพราะเหล็กเปลี่ยนสถานะได้
\end{enumerate}
\item หน่วยความดัน Pascal (Pa) มีค่าเท่ากับข้อใด \hfill \textbf{\small (Main)}
\begin{enumerate}[label=)]
\item ก) 1 N/m
\item ข) 1 N/m²
\item ค) 1 kg/m³
\item ง) 1 J/m
\end{enumerate}
\item เครื่องอัดไฮดรอลิกมีลูกสูบเล็กพื้นที่ 0.01 m² กดด้วยแรง 100 N และลูกสูบใหญ่พื้นที่ 0.1 m² จะยกวัตถุได้กี่นิวตัน \hfill \textbf{\small (Main)}
\begin{enumerate}[label=)]
\item ก) 100 N
\item ข) 1,000 N
\item ค) 10,000 N
\item ง) 10 N
\end{enumerate}
\item น้ำไหลในท่อระดับพื้นที่ 0.02 m² ด้วยความเร็ว 3 m/s และความดัน 150,000 Pa ที่จุดที่ท่อหดเหลือ 0.01 m² ความดันเป็นเท่าไร (ρ น้ำ = 1,000 kg/m³) \hfill \textbf{\small (Main)}
\begin{enumerate}[label=)]
\item ก) 75,000 Pa
\item ข) 112,500 Pa
\item ค) 150,000 Pa
\item ง) 187,500 Pa
\end{enumerate}
\item ความดันที่ก้นภาชนะที่บรรจุของเหลวสูง 50 cm ขึ้นกับปริมาตรของภาชนะหรือไม่ \hfill \textbf{\small (Main)}
\begin{enumerate}[label=)]
\item ก) ขึ้นกับ เพราะถ้าภาชนะกว้างขึ้นจะรองรับน้ำหนักมากขึ้น
\item ข) ไม่ขึ้น เพราะความดันขึ้นกับความลึกเท่านั้น
\item ค) ขึ้นกับรูปร่างของภาชนะ
\item ง) ขึ้นกับมวลของของเหลว
\end{enumerate}
\item วัตถุมีน้ำหนักในอากาศ 50 N จมในน้ำได้ 30 N แรงลอยมีค่าเท่าไร \hfill \textbf{\small (Main)}
\begin{enumerate}[label=)]
\item ก) 20 N
\item ข) 30 N
\item ค) 50 N
\item ง) 80 N
\end{enumerate}
\item น้ำไหลในท่อเส้นผ่านศูนย์กลาง 4 cm ด้วยอัตราการไหล 0.02 m³/s ความเร็วเฉลี่ยของน้ำเป็นเท่าไร \hfill \textbf{\small (Main)}
\begin{enumerate}[label=)]
\item ก) 0.64 m/s
\item ข) 1.6 m/s
\item ค) 6.4 m/s
\item ง) 16 m/s
\end{enumerate}
\item สสารสถานะใดมีความหนาแน่นมากที่สุดโดยทั่วไป \hfill \textbf{\small (Main)}
\begin{enumerate}[label=)]
\item ก) ของแข็ง
\item ข) ของเหลว
\item ค) แก๊ส
\item ง) ขึ้นกับสถานะการณ์
\end{enumerate}
\item น้ำไหลในท่อระดับที่ความดัน 200 kPa และความเร็ว 3 m/s ที่จุดที่ความดันลดเหลือ 150 kPa ความเร็วจะเป็นเท่าไร \hfill \textbf{\small (Main)}
\begin{enumerate}[label=)]
\item ก) 1.5 m/s
\item ข) √5 m/s
\item ค) ≈ 6 m/s
\item ง) 4.5 m/s
\end{enumerate}
\item ข้อใดเป็นการประยุกต์ใช้หลักของปาสคาล \hfill \textbf{\small (Main)}
\begin{enumerate}[label=)]
\item ก) เรือลอยน้ำ
\item ข) เครื่องอัดไฮดรอลิก
\item ค) เครื่องบินบินได้
\item ง) ลูกโป่งลอยขึ้น
\end{enumerate}
\item ที่ความลึก 10 m ในน้ำ ความดันเกจ (เมื่อเทียบกับบรรยากาศ) มีค่าเท่าไร (g = 10 m/s²) \hfill \textbf{\small (Main)}
\begin{enumerate}[label=)]
\item ก) 10,000 Pa
\item ข) 100,000 Pa
\item ค) 200,000 Pa
\item ง) 1,000,000 Pa
\end{enumerate}
\item ลูกบอลพลาสติกรัศมี 10 cm ลอยในน้ำโดยจมครึ่งหนึ่ง มวลของลูกบอลเป็นเท่าไร (ρ น้ำ = 1,000 kg/m³, g = 10 m/s²) \hfill \textbf{\small (Main)}
\begin{enumerate}[label=)]
\item ก) ≈ 2.1 kg
\item ข) ≈ 4.2 kg
\item ค) ≈ 0.5 kg
\item ง) ≈ 1 kg
\end{enumerate}
\item น้ำไหลในท่อพื้นที่ 0.01 m² ด้วยความเร็ว 5 m/s ไปยังท่อเล็กพื้นที่ 0.0025 m² อัตราการไหลปริมาตรเป็นเท่าไร \hfill \textbf{\small (Main)}
\begin{enumerate}[label=)]
\item ก) 0.0125 m³/s
\item ข) 0.05 m³/s
\item ค) 0.2 m³/s
\item ง) 0.005 m³/s
\end{enumerate}
\item น้ำ 1 ลิตร (0.001 m³) มีมวลเท่าไร \hfill \textbf{\small (Extra)}
\begin{enumerate}[label=)]
\item ก) 1 g
\item ข) 100 g
\item ค) 1 kg
\item ง) 10 kg
\end{enumerate}
\item ถ้าแรง 100 N กระทำตั้งฉากกับพื้นที่ 2 m² ความดันมีค่าเท่าไร \hfill \textbf{\small (Extra)}
\begin{enumerate}[label=)]
\item ก) 50 Pa
\item ข) 200 Pa
\item ค) 0.02 Pa
\item ง) 500 Pa
\end{enumerate}
\item ถ้อยคำใดต่อไปนี้ถูกต้องเกี่ยวกับความดันในของเหลว \hfill \textbf{\small (Extra)}
\begin{enumerate}[label=)]
\item ก) ความดันขึ้นกับทิศทาง
\item ข) ความดันที่ความลึกเท่ากันในของเหลวเดียวกันจะเท่ากันทุกทิศทาง
\item ค) ความดันขึ้นกับปริมาตรของของเหลว
\item ง) ความดันที่ก้นภาชนะขึ้นกับรูปร่างภาชนะ
\end{enumerate}
\item วัตถุลอยในน้ำ โดย 30\% จมอยู่ใต้น้ำ ถ้าน้ำหนักวัตถุ = 30 N แรงลอยมีค่าเท่าไร \hfill \textbf{\small (Extra)}
\begin{enumerate}[label=)]
\item ก) 3 N
\item ข) 9 N
\item ค) 30 N
\item ง) 100 N
\end{enumerate}
\item ตามหลักของเบอร์นูลลี สมการ P + ½ρv² + ρgh = คงที่ เรียกอะไร \hfill \textbf{\small (Extra)}
\begin{enumerate}[label=)]
\item ก) สมการความต่อเนื่อง
\item ข) สมการเบอร์นูลลี
\item ค) สมการของปาสคาล
\item ง) สมการของอาร์คิมีดีส
\end{enumerate}
\item อัตราการไหลปริมาตร (Q) มีหน่วยเป็นอะไรในระบบ SI \hfill \textbf{\small (Extra)}
\begin{enumerate}[label=)]
\item ก) m/s
\item ข) m²/s
\item ค) m³/s
\item ง) kg/s
\end{enumerate}
\item เครื่องอัดไฮดรอลิกมีอัตราส่วนพื้นที่ลูกสูบ 10:1 ถ้าใช้แรง 50 N กดลูกสูบเล็ก จะได้แรงขนาดเท่าใดที่ลูกสูบใหญ่ \hfill \textbf{\small (Extra)}
\begin{enumerate}[label=)]
\item ก) 500 N
\item ข) 5 N
\item ค) 5000 N
\item ง) 250 N
\end{enumerate}
\item ความดันบรรยากาศปกติ 1 atm = 1.013×10⁵ Pa ถ้านักดำน้ำลงไป 30 m ความดันรวมที่เขาแบกรับเป็นเท่าไร \hfill \textbf{\small (Extra)}
\begin{enumerate}[label=)]
\item ก) 1.013×10⁵ Pa
\item ข) 2.013×10⁵ Pa
\item ค) 3.013×10⁵ Pa
\item ง) 4.013×10⁵ Pa
\end{enumerate}
\item ทำไมลูกโป่งลมร้อนจึงลอยขึ้น \hfill \textbf{\small (Extra)}
\begin{enumerate}[label=)]
\item ก) เพราะอากาศร้อนเบากว่าอากาศเย็น
\item ข) เพราะลูกโป่งมีความหนาแน่นน้อยกว่าอากาศเย็นรอบข้าง
\item ค) เพราะอากาศร้อนขยายตัว
\item ง) ถูกทั้ง ก และ ข
\end{enumerate}
\item เครื่องบินบินได้เพราะเหตุผลใด \hfill \textbf{\small (Extra)}
\begin{enumerate}[label=)]
\item ก) เครื่องยนต์ผลักให้ไปข้างหน้า
\item ข) รูปร่างปีกทำให้ความดันอากาศด้านบนปีกต่ำกว่าด้านล่าง สร้างแรงยก
\item ค) เครื่องบินเบากว่าอากาศ
\item ง) แรงลอยจากอากาศร้อน
\end{enumerate}
\item วัตถุ A มีความหนาแน่น 0.8 g/cm³ และวัตถุ B มีความหนาแน่น 1.2 g/cm³ ถ้าทั้งสองลอยน้ำ วัตถุใดจมมากกว่า \hfill \textbf{\small (Extra)}
\begin{enumerate}[label=)]
\item ก) วัตถุ A
\item ข) วัตถุ B
\item ค) จมเท่ากัน
\item ง) ไม่สามารถระบุได้
\end{enumerate}
\item น้ำไหลในท่อที่มีเส้นผ่านศูนย์กลาง 10 cm ด้วยความเร็ว 2 m/s ถ้าท่อแคบลงเหลือเส้นผ่านศูนย์กลาง 5 cm ความเร็วจะเป็นเท่าไร \hfill \textbf{\small (Extra)}
\begin{enumerate}[label=)]
\item ก) 1 m/s
\item ข) 4 m/s
\item ค) 8 m/s
\item ง) 16 m/s
\end{enumerate}
\item วัตถุมวล 5 kg ปริมาตร 0.006 m³ ถูกปล่อยลงในน้ำ (ρ = 1,000 kg/m³) วัตถุจะลอยหรือจม \hfill \textbf{\small (Extra)}
\begin{enumerate}[label=)]
\item ก) ลอย เพราะ F\_b > W
\item ข) ลอย เพราะ F\_b = W
\item ค) จม เพราะ F\_b < W
\item ง) จม เพราะ F\_b = W
\end{enumerate}
\item 1 atm มีค่าเท่ากับกี่ mmHg \hfill \textbf{\small (Extra)}
\begin{enumerate}[label=)]
\item ก) 760 mmHg
\item ข) 100 mmHg
\item ค) 10 mmHg
\item ง) 1,000 mmHg
\end{enumerate}
\item น้ำไหลในท่อระดับพื้นที่ 0.01 m² ความเร็ว 4 m/s ความดัน 100 kPa ที่ท่อหดพื้นที่ 0.004 m² ความดันเป็นเท่าไร \hfill \textbf{\small (Extra)}
\begin{enumerate}[label=)]
\item ก) 75 kPa
\item ข) 100 kPa
\item ค) 125 kPa
\item ง) 150 kPa
\end{enumerate}
\item ภาชนะสองใบรูปต่างกันมีน้ำระดับเท่ากันที่ 30 cm จากก้น ความดันที่ก้นทั้งสองภาชนะเปรียบเทียบอย่างไร \hfill \textbf{\small (Extra)}
\begin{enumerate}[label=)]
\item ก) ภาชนะกว้างกว่ามีความดันมากกว่า
\item ข) ภาชนะแคบกว่ามีความดันมากกว่า
\item ค) ความดันเท่ากัน
\item ง) ขึ้นกับปริมาตรน้ำ
\end{enumerate}
\item เรือลอยน้ำเค็ม (ความหนาแน่น 1,025 kg/m³) กับน้ำจืด (1,000 kg/m³) ที่ไหนเรือจมลึกกว่า \hfill \textbf{\small (Extra)}
\begin{enumerate}[label=)]
\item ก) น้ำเค็ม เพราะแรงลอยน้อยกว่า
\item ข) น้ำจืด เพราะแรงลอยน้อยกว่า
\item ค) เท่ากัน เพราะน้ำหนักเรือเท่าเดิม
\item ง) น้ำจืด เพราะความหนาแน่นน้อยกว่า
\end{enumerate}
\item ท่อน้ำพื้นที่ 50 cm² มีน้ำไหล 10 L/s ความเร็วเป็นเท่าไร (1 L = 1,000 cm³) \hfill \textbf{\small (Extra)}
\begin{enumerate}[label=)]
\item ก) 0.2 m/s
\item ข) 2 m/s
\item ค) 5 m/s
\item ง) 20 m/s
\end{enumerate}
\item กระบอกฉีดน้ำทำงานบนหลักใด \hfill \textbf{\small (Extra)}
\begin{enumerate}[label=)]
\item ก) หลักของอาร์คิมีดีส
\item ข) หลักของปาสคาล
\item ค) หลักของเบอร์นูลลี
\item ง) หลักของชาร์ลส์
\end{enumerate}
\item ข้อใดไม่ใช่ตัวแปรในสูตร P = ρgh \hfill \textbf{\small (Extra)}
\begin{enumerate}[label=)]
\item ก) ρ คือความหนาแน่นของของเหลว
\item ข) g คือความเร่งโน้มถ่วง
\item ค) h คือความสูงของภาชนะ
\item ง) P คือความดันเกจ
\end{enumerate}
\item ทรงกลมพลาสติกรัศมี 5 cm มีมวล 0.5 kg ถูกปล่อยลงในน้ำ ทรงกลมจะ \hfill \textbf{\small (Extra)}
\begin{enumerate}[label=)]
\item ก) ลอยเฉยๆ
\item ข) ลอยโดยจมไป 80\%
\item ค) จมลงไปทั้งหมด
\item ง) ลอยโดยจมไป 50\%
\end{enumerate}
\item ถ้าอัตราการไหลปริมาตรของน้ำในท่อเป็น 0.03 m³/s และท่อมีพื้นที่หน้าตัด 0.015 m² ความเร็วน้ำเป็นเท่าไร \hfill \textbf{\small (Extra)}
\begin{enumerate}[label=)]
\item ก) 0.45 m/s
\item ข) 2 m/s
\item ค) 4.5 m/s
\item ง) 20 m/s
\end{enumerate}
\item ทำไมน้ำจึงพุ่งออกจากสายน้ำแบบพลาสติกได้ไกลกว่าเมื่อบีบแน่น \hfill \textbf{\small (Extra)}
\begin{enumerate}[label=)]
\item ก) เพราะความดันอากาศภายนอกลดลง
\item ข) เพราะความเร็วน้ำเพิ่มขึ้นเมื่อพื้นที่หน้าตัดเล็กลง
\item ค) เพราะน้ำขยายตัว
\item ง) เพราะแรงโน้มถ่วง
\end{enumerate}
\item ความดัน 2 atm มีค่าเท่ากับกี่ Pa \hfill \textbf{\small (Extra)}
\begin{enumerate}[label=)]
\item ก) 2×10⁵ Pa
\item ข) 2.026×10⁵ Pa
\item ค) 1.013×10⁵ Pa
\item ง) 4.052×10⁵ Pa
\end{enumerate}
\item ถ้าวัตถุลอยในของเหลว A (ความหนาแน่น 800 kg/m³) โดยจม 60\% และลอยในของเหลว B โดยจม 80\% ของเหลวใดมีความหนาแน่นมากกว่า \hfill \textbf{\small (Extra)}
\begin{enumerate}[label=)]
\item ก) ของเหลว A
\item ข) ของเหลว B
\item ค) เท่ากัน
\item ง) ไม่สามารถระบุได้
\end{enumerate}
\item ข้อใดกล่าวถูกต้องเกี่ยวกับความดันในของเหลว \hfill \textbf{\small (Extra)}
\begin{enumerate}[label=)]
\item ก) ความดันที่ผิวบนเท่ากับความดันที่ก้นภาชนะ
\item ข) ความดันเพิ่มขึ้นเมื่อลึกลงไป
\item ค) ความดันที่ด้านข้างภาชนะน้อยกว่าความดันที่ก้น
\item ง) ข้อ ข และ ค ถูก
\end{enumerate}
\item ในสมการเบอร์นูลลี P₁ + ½ρv₁² + ρgh₁ = P₂ + ½ρv₂² + ρgh₂ ถ้าท่ออยู่ในแนวราบ (h₁ = h₂) และความเร็วเพิ่มขึ้น 3 เท่า ความดันจะเปลี่ยนอย่างไร \hfill \textbf{\small (Extra)}
\begin{enumerate}[label=)]
\item ก) เพิ่มขึ้น 9 เท่า
\item ข) ลดลง 4 เท่า
\item ค) ลดลง 4.5 เท่าจากค่าเดิม
\item ง) เพิ่มขึ้น 4 เท่า
\end{enumerate}
\item สารมวล 500 g มีปริมาตร 250 cm³ สารนี้จะลอยหรือจมในน้ำ \hfill \textbf{\small (Extra)}
\begin{enumerate}[label=)]
\item ก) ลอย เพราะความหนาแน่น 2 g/cm³ > 1
\item ข) ลอย เพราะความหนาแน่น 2 g/cm³ < 1
\item ค) จม เพราะความหนาแน่น 2 g/cm³ > 1
\item ง) จม เพราะความหนาแน่น 2 g/cm³ < 1
\end{enumerate}
\item ก้อนน้ำแข็งลอยในน้ำ ถ้าน้ำแข็งละลาย ระดับน้ำในภาชนะจะเป็นอย่างไร (ไม่คิดการระเหย) \hfill \textbf{\small (Extra)}
\begin{enumerate}[label=)]
\item ก) สูงขึ้น
\item ข) ต่ำลง
\item ค) เท่าเดิม
\item ง) ขึ้นกับปริมาตรน้ำแข็ง
\end{enumerate}
\item ท่อพลาสติกมีน้ำไหลผ่าน ที่จุด A มีพื้นที่ใหญ่และความดันสูง ที่จุด B มีพื้นที่เล็กและความดันต่ำ ข้อใดอธิบายได้ถูกต้อง \hfill \textbf{\small (Extra)}
\begin{enumerate}[label=)]
\item ก) P สูง → A ใหญ่ → v ต่ำ และ P ต่ำ → A เล็ก → v สูง ตามสมการความต่อเนื่อง
\item ข) P สูง → v สูง เสมอ
\item ค) A ใหญ่ → v สูง เสมอ
\item ง) P และ v ไม่มีความสัมพันธ์กัน
\end{enumerate}
\end{enumerate}
\end{document}
